% Created 2026-09-05 Sat 16:29
% Intended LaTeX compiler: xelatex
\documentclass[12pt]{article}
\usepackage[utf8]{inputenc}
\usepackage[T1]{fontenc}
\usepackage{graphicx}
\usepackage{longtable}
\usepackage{wrapfig}
\usepackage[margin=1in]{geometry}
\usepackage{rotating}
\usepackage{setspace}
\onehalfspacing
\usepackage[normalem]{ulem}
\usepackage{amsmath}
\usepackage{parskip}
\usepackage{fontspec}
\setmainfont{Libertinus Serif}
\usepackage{amssymb}
\usepackage{capt-of}
\setcounter{secnumdepth}{4} % Tells LaTeX to number down to paragraphs
\setcounter{tocdepth}{4}  
\usepackage{hyperref}
\usepackage[style=oscola, backend=biber]{biblatex}
\addbibresource{/Users/krishnajani/Documents/Libib.bib}
\usepackage{csquotes}
\usepackage[british]{babel}
\setlength{\parskip}{0.5\baselineskip}
\author{Krishna Jani}
\date{\today}
\title{Orator Marketing v. Union of India}

\begin{document}
\maketitle

\section*{Facts}

A short question was involved in this appeal, that is, whether a loan given to a corporate person \textbf{free of interest} on account of working capital requirements can be characterised as a financial debt for the purposes of Section 7 of the IBC.

\begin{itemize}
\item In this case, S advanced a loan of 1.6 crores to CD for a period of 2 years. The loan was due to be repaid in a particular period, of which a substantial sum remains outstanding. 
\item In furtherance of this, the appellant filed a petition under Section 7 for initiation of CIRP process. 
\item The petition was rejected by the NCLT on the grounds that the applicant has failed to prove that the loan was disbursed against consideration for the \textbf{time value of mone}y (particularly when the respondent company has affirmed that no interest has been paid or payable at that or at any point in time).
\item Being aggrieved, an appeal was further filed before the NCLAT, which upheld the rejection. The impugned order states that: Only that money which is borrowed against the payment of interest constitutes financial debt. \textbf{If money borrowed is not against payment of interest, the core requirement as to time value of money does not exist,} and thus it does not constitute a financial debt, and the Section 7 application is incompetent. 
 
\end{itemize}

\section*{Court's Position}

The Court held that the judgement and order of the NCLAT is patently flawed. Both the NCLAT and NCLT have misconstrued the definition of financial debt found in Section 5(8) of the IBC by reading the same in isolation and out of context. Section 5(8) defines financial debt to mean a debt, along with interest, if any, which is disbursed against the consideration of time value of money, and includes money borrowed against payment of interest as per Section 5(8A) of the IBC. The NCLT and NCLAT have overlooked the words "if any", which could not have been intended to be otiose. If there is no interest payable on the loan, only the outstanding principal would qualify as a financial debt.

The tribunals have also failed to notice that the residual clause in Section 5(8)(f) includes any amounts raised under any transaction having the commercial effect of borrowing. Subclauses (a) to (i) of Section 5(8) are illustrative and are not exhaustive.

In Swiss Ribbons\footcite{justicer.fnarimanSwissRibbonsPvt2019} it was held the Court identified the substantive position of an FC as follows:

\begin{quote}
  A financial creditor is a person who has direct engagement in the functioning of the corporate debtor, who is involved right from the beginning while assessing the viability of the corporate debtor, who would engage in restructuring of the loan as well as reorganisation of the corporate debtor's business when there is financial stress. Therefore, the financial creditor acquires a unique position and could be entrusted with the task of ensuring the sustenance and growth of the corporate debtor according to that of a guardian. 
\end{quote}

In sum and substance, the financial creditor becomes an important stakeholder in the financial health of the company, and this position must be taken into consideration rather than fixating on the presence or the absence of interest (since the statute itself states that the absence of interest shall not render the debt outside the paradigm of financial debt). 


The court alluded to the golden rule, which requires that the words of a statute must be \textit{prima facie} given, their ordinary meaning when the language or the phrasing employed by the statute is precise and plain.


It is noted that the definition of ``debt'' is also inclusive and the definition of ``financial debt'' does not expressly exclude interest-free loans. Through its plain reading, it would mean to include interest-free loans advanced to for business operations would constitute financial debt.

The order of the NCLAT was therefore set aside, and the appeal was allowed. 


\end{document}
